\newpage
\section{Discussion and Conclusion}
\label{sec:discussion}

\subsection{Discussion}

\paragraph{Comparing with the hypothesis.}
Our results agree with all three parts of the hypothesis
(Section~\ref{sec:hypothesis}). In the baseline, a 1~g fragment changes
the satellite's velocity by $20.0\unit{mm\,s^{-1}}$, which is only
$2.6\times10^{-6}$ of the orbital speed (Table~\ref{tab:results}). So one
hit like this does not really change the orbit. The spin is
$1.65^\circ\unit{s^{-1}}$, about what we expected. The angular momentum
$m v_0\cos\theta\, b=34.6\unit{kg\,m^2\,s^{-1}}$ is large because the hit
is at the tip of the panel, 4~m from the CoM. At this rate the satellite
would turn a full circle in less than four minutes if nothing stopped it.
Finally, $99.9988\%$ of the $50\unit{kJ}$ is lost. This is a little less
than the $M/(m+M)=99.9998\%$ for a hit through the CoM, because
$0.5\unit{J}$ goes into rotation. The satellite gets less than $1\unit{J}$
of kinetic energy in total. In other words, the momentum is passed on
very well, but the energy is not.

\paragraph{What the sweeps show.}
The angle sweep (Fig.~\ref{fig:angle}) shows that the type of collision
changes not only how much momentum is passed on, but also which part. A
fragment that sticks always gives all of its momentum, so $|\Delta\vect V|$
does not depend on $\theta$. A fragment that bounces off only doubles the
normal part, so it gives a bigger push than a sticking fragment only when
$\theta<60^\circ$. The spin depends only on the normal impulse in our
geometry, so it drops like $\cos\theta$ in both cases. The mass--speed map
(Fig.~\ref{fig:mass-speed}) grows linearly with $m$ and $v_0$ as long as
$m\ll M$. Its top corner (kilograms at $15\unit{km\,s^{-1}}$, up to
$1.1\unit{GJ}$) is outside what our model can describe: a hit like that
would break the satellite apart instead of just pushing it.

\paragraph{How precise are the inputs?}
When $m\ll M$, $|\Delta\vect V|$ and $\omega$ are proportional to $m$ and
$v_0$. So a 10\% error in the input mass or speed gives a 10\% error in
the results, and $E_{\mathrm{diss}}$ doubles any relative error in $v_0$.
The computer's precision is not a problem: the errors of $10^{-16}$
(Table~\ref{tab:results}) and the match with the 1-D formulas to
$6.7\times10^{-16}$ show that the code solves the equations exactly. What
limits the accuracy is the model itself. The satellite is rigid and 2-D,
the fragment is a point, and a real hit at these speeds is somewhere
between elastic and completely inelastic. Material thrown back out of the
crater can even make the momentum transfer bigger than in the completely
inelastic case, and our model does not include this.

\paragraph{Comparing with a real event.}
For the Sentinel-1A hit, \citet{krag2017} used GPS and attitude data to
estimate a velocity change of only about $0.7\unit{mm\,s^{-1}}$ and a
peak rotation rate of about $0.04^\circ\unit{s^{-1}}$ about the yaw axis.
Our numbers are bigger because our satellite is lighter, with a smaller
moment of inertia, and because we chose the fragment and its speed
ourselves. But the overall picture is the same: the velocity change was
tiny compared with the orbital speed, while the change in orientation
could clearly be measured. This supports our conclusion that, for small
debris, the spin from an off-centre hit matters more than the change in
orbit.

\subsection{Applications}
\label{sec:applications}

The same momentum and energy calculations are useful in real engineering.
\begin{itemize}
  \item \textit{Attitude control.} The spin from
    Eq.~\eqref{eq:omega-inelastic} is a disturbance that the reaction
    wheels or thrusters have to cancel. The angular momentum from a hit,
    $m v_n b$, tells engineers how much the wheels must be able to absorb.
    For example, Sentinel-1A had a sudden change of orientation after its
    2016 hit \citep{krag2017}.
  \item \textit{Keeping the orbit.} Comparing $|\Delta\vect V|$ with
    $v_{\mathrm{orb}}$ tells us whether a hit needs an orbit correction.
    For debris from millimetres to centimetres in size the change is tiny,
    but a fragment of several kilograms would need a correction burn.
  \item \textit{Shielding and panel design.} Equation~\eqref{eq:ediss}
    shows that in an inelastic hit almost all the energy is released at
    one spot. This is why spacecraft use spaced shields (Whipple shields)
    that break the fragment apart before it reaches the main wall, and why
    solar panels are built so that one crater only knocks out part of the
    power.
  \item \textit{The debris problem.} Every hit can create new fragments.
    Working out the energy and momentum of collisions is the first step in
    models of the debris population and of the chain reaction predicted by
    \citet{kessler1978}.
\end{itemize}


\subsection{Conclusion}

Using conservation of linear and angular momentum together with the
impulse--momentum theorem and the parallel-axis theorem, we could solve
exactly what happens when a debris fragment hits a solar panel, for both
the elastic and the completely inelastic case. Our Python program shows
that a 1~g fragment at $10\unit{km\,s^{-1}}$ changes the velocity of a
500~kg satellite by only $20\unit{mm\,s^{-1}}$ ($2.6\times10^{-6}$ of its
orbital speed). However, it makes the satellite spin at
$1.65^\circ\unit{s^{-1}}$, and almost all of the $50\unit{kJ}$ impact
energy is lost at the panel. These results confirm our hypothesis, show
the same pattern as the Sentinel-1A measurements, and pass the
conservation checks to round-off level. As a next step, the two ideal
cases could be replaced by a measured momentum-transfer value, and the
orbit after the hit could be calculated.
